% Intended LaTeX compiler: xelatex
\documentclass[14pt,aspectratio=169]{beamer}
\usepackage{graphicx}
\usepackage{longtable}
\usepackage{wrapfig}
\usepackage{rotating}
\usepackage[normalem]{ulem}
\usepackage{amsmath}
\usepackage{amssymb}
\usepackage{capt-of}
\usepackage{hyperref}
\author{i}
\date{\today}
\title{The Pittsburgh template}
\subtitle{A beamer theme based on metropolis}
\author{Ian Waudby-Smith}
\date{}
\institute{Institute\\Pittsburgh, PA, USA}
\newcommand{\includetoc}{1} % set 1 or 0 to include Table of Contents or not
\newcommand{\marginsize}{2cm} % Margin size globally
\newcommand{\additionaltitlemargin}{0cm} % Additional margin for title page
\definecolor{accent}{RGB}{192 97 106} % Accent color throughout
\usepackage{../../beamertheme-pittsburgh/beamerthemepittsburgh}
\usepackage{../../beamertheme-pittsburgh/pittsburgh_fancyfonts}
\usefonttheme{serif} % default family is serif
% \setmainfont{Palatino} % Nice font but commented out because it's not available on some OSs by default

\begin{document}

\newcommand{\Accent}[1]{\textbf{\textcolor{accent}{#1}}}

\maketitle
%%%%%%%%%%%%%%%%%%%%%%%%%%%%%%%%
% Collaborators slide
%%%%%%%%%%%%%%%%%%%%%%%%%%%%%%%%

\begin{frame}[noframenumbering]{Collaborators}
\thispagestyle{empty}

\newcommand{\collaborators}{
  {carnegie.jpeg/Collaborator 1\\Carnegie Tech},
  {mellon.jpg/Collaborator 2\\Mellon Institute}}

% First argument (./assets) is the location of images,
% Second argument is the collaborators list
% Third argument is the width of minipages
% Fourth argument is the height of the images
% Fifth argument is the size of text
\displayimages{./assets/}{\collaborators}{0.45}{0.35}{\small}

\end{frame}

\renewcommand{\toctitle}{Outline of demo presentation}
\ifnum\includetoc=1{%
    \firstpagetoc
  }
\fi

\section{An overview of the \Accent{pittsburgh} template}
\begin{frame}{What is the \Accent{pittsburgh} template?}
 \Accent{pittsburgh} is a beamer template based on the widely used \texttt{metropolis} theme.

It's called \Accent{pittsburgh} because Pittsburgh, PA is the metropolis where this was made.
\end{frame}

\begin{frame}{Why?}
  \small
 I love the aesthetics of Apple's Keynote presentations on macOS but
\begin{enumerate}
\item Keynote only runs on macOS and thus presentations are lost when switching OSs,
\pause
\item Its files tend to be quite large relative to their content, and
\pause
\item Its math fonts don't match the \LaTeX{} ones.
\pause
\end{enumerate}
Beamer fixes the above and \Accent{pittsburgh} attempts to carry over some of the slick design of Keynote.
\end{frame}


\section{Some cool features}

\begin{frame}{Cool feature I: \emph{Pretty text boxes}}
\widen[1cm]{
  \begin{pghcolorbox}{turquoise}{\textbf{Contextual bandits} (Langford \& Zhang, '07)}
  

    For $t= 1,2,...$:

    \quad Agent observes a context $X_t $ (typically in $\mathbb R^d$).

    \quad Agent takes action $A_t \sim h_t(\cdot \mid X_t)$.

    \quad Agent observes reward $R_t$ (typically in $\mathbb R$).

    EndFor

  \end{pghcolorbox}
}
\end{frame}

\begin{frame}{Cool feature II: \emph{Adjustable slide widths}}
 What if a particular sentence just barely fits in one line?

(Like the one above.) The \texttt{\textbackslash widen} command allows you to dynamically widen or shrink the allowable width:

\widen[1.5cm]{\centering ``What if a particular sentence just barely fits in one line?''}
\end{frame}

\begin{frame}{}
\centering
Have fun!
\end{frame}


\appendixpgh


\begin{frame}{Appendix slide}
  Start the appendix with the \texttt{\textbackslash appendixpgh} command. Slide numbers will automatically enumerate as ``A1'', ``A2'', and so on.
\end{frame}

\end{document}


%%% Local Variables: 
%%% coding: utf-8
%%% mode: latex
%%% TeX-engine: xetex
%%% End: 
